\begin{table}[H]
\small
\setlength{\tabcolsep}{4.2pt}
\renewcommand{\arraystretch}{0.95}

\caption{\label{tab:ea20-income-gap-product-contributions}Product-group contributions to euro-area cumulative income inflation gaps, June 2021--June 2023}
\centering
\begin{tabular}[t]{lllll}
\toprule
Product group & Cumulative inflation & Q1 contribution & Q5 contribution & Inflation gap\\
\midrule
Food and non-alcoholic beverages & 4.3 & 4.5 & 4.2 & 0.3\\
Alcoholic beverage, tobacco and narcotics & 0.5 & 0.7 & 0.5 & 0.2\\
Clothing and footwear & 0.3 & 0.4 & 0.6 & -0.2\\
Housing (rentals and repairs) & 1.2 & 0.7 & 0.8 & -0.1\\
Water, electricity, gas and other fuels & 1.5 & 6.4 & 5.6 & 0.8\\
\addlinespace
Transport & 1.9 & 1.6 & 3.0 & -1.4\\
Other & 5.0 & 0.3 & 0.5 & -0.2\\
\midrule
\textbf{Total} & \textbf{14.7} & \textbf{14.6} & \textbf{15.2} & \textbf{-0.6}\\
\bottomrule
\end{tabular}
\vspace{-0.4em}
\begin{flushleft}
\footnotesize{Notes. Entries are percentage points. Product-group contributions are cumulated from the month after June 2021 through June 2023  and rescaled so that the total row reproduces the euro-area cumulative price-index changes reported in Table~\ref{tab:country-heterogeneity-cumulative-gap}. Product groups use the default mapping in \texttt{plot\_contribution\_gap}. France is computed at COICOP level 3 with INSEE income quintiles; the euro area is rebuilt from country-level contributions and EA20 HICP country weights. The gap is Q1 minus Q5. Sources: Eurostat, authors' calculations.}
\end{flushleft}
\end{table}
